This short section restates the one geometric definition on which the zero
theorems of Section~\ref{sec:zero-geometry} depend.  It is imported from Paper~0,
where the definition, the frame equivalence, the models, the flow, the
commutativization, and the contact structure are developed and proved; the
statements below are repeated here only so that this paper can be read on its
own, and they carry the same labels and hypotheses as their source.

\begin{definition}[Regular arithmetic expression space, imported]
\label{def:regular-aes}
A \emph{regular arithmetic expression space}, or regular AES, is a tuple
\[
 \mathfrak E=(M,g,a;\mu,\lambda)
\]
such that $M$ is an oriented smooth surface, possibly with boundary, $g$ is a
smooth Riemannian metric, $a\in C^\infty(M,\mathbb R)$,
$\mu\in\mathbb R\setminus\{0\}$, $\lambda\in\mathbb R$, and
\begin{equation}\label{eq:regular-aes-eikonal}
 |\nabla a|_g^2=\mu^2+\lambda^2a^2
\end{equation}
everywhere on $M$.  The function $a$ is called the \emph{assignment function}.
\end{definition}

The metric is part of the structure; \eqref{eq:regular-aes-eikonal} is an
eikonal equation on a specified Riemannian surface, not a metric-free identity.
Paper~0 defines the arithmetic motions of such a space and proves that they
realize addition, subtraction, multiplication, and division by a constant as
motions of $M$.

\begin{proposition}[Canonical arithmetic frame, imported]
\label{prop:canonical-arithmetic-frame}
For an oriented Riemannian surface $(M,g)$, a smooth function $a$, and constants
$\mu\ne0$ and $\lambda$, the following are equivalent.
\begin{enumerate}
\item The eikonal identity \eqref{eq:regular-aes-eikonal} holds.
\item There is a unique positively oriented orthonormal frame $(X_u,X_v)$
      satisfying
      \begin{equation}\label{eq:arithmetic-frame-derivatives}
       X_u a=\mu,
       \qquad
       X_v a=\lambda a.
      \end{equation}
\end{enumerate}
When these conditions hold, the frame is
\begin{equation}\label{eq:explicit-arithmetic-frame}
 X_u=\frac{\mu}{q}E_a-\frac{\lambda a}{q}JE_a,
 \qquad
 X_v=\frac{\lambda a}{q}E_a+\frac{\mu}{q}JE_a,
 \qquad
 E_a=\frac{\nabla a}{q},
 \qquad
 q=\sqrt{\mu^2+\lambda^2a^2},
\end{equation}
and it satisfies
\begin{equation}\label{eq:arithmetic-frame-bracket-on-a}
 [X_u,X_v]a=\mu\lambda .
\end{equation}
\end{proposition}

The proof is given in Paper~0.  We use only the following immediate consequences,
which is why the statement is repeated rather than cited.

\begin{enumerate}
  \item $\mu\ne0$ makes $q$ everywhere positive, so $|\nabla a|$ never vanishes.
  This is the source of the regularity of zero sets below.
  \item The bracket identity \eqref{eq:arithmetic-frame-bracket-on-a} shows that,
  when $\lambda\ne0$, the two arithmetic directions cannot be realized as a
  coordinate frame.
\end{enumerate}

No other object of Paper~0 is repeated in this paper.  In particular, the basic
hyperbolic model $\mathfrak E_0$, the punctured model $\mathfrak E_1$, the
continuous arithmetic motion, the affine cocycles, the accumulative commutative
space, tearing, and the contact structure are cited where needed and never
reproved here.
